\subsubsection{LINGUAGEM DE PROGRAMAÇÃO: \en{GOLANG}}

\en{Golang} é uma linguagem de programação simimilar a linguagem C, dela herdou a sintáse de expressões, tipos de dados básicos e a ênfase em programas que compilam para código de máquina eficiente. Criada por Robert Griesemer, Rob Pike, e Ken Thompson, todos da \en{Google}, em 2007 para resolver o problema crescente da complexidade dos sistema da empresa. \en{Go} é apropriado para a contrução de servidores que atuam como infraestutura em rede \cite{donovan2016gopl}, portanto a escolha para a construção do \en{middleware}. 

A linguagem coloca muita ênfase na eficiência, \en{Go} possui funcionalidades inspiradas em outras linguagens, como: \en{garbage collection} de linguagens como \en{Java}, o conceito de pacotes de \en{Modula-2}, a sintaxe de métodos de \en{Oberon-2}, dentre outras funcionalidades. As \en{goroutines}, fluxos leves de execução, são outro aspecto da língua que demonstra a preocupação por eficiência, onde: criar uma \en{goroutine} é barato e criar milhões é prático \cite{donovan2016gopl}. 

Alguns exemplos de produtos que usam \en{Go} são: \en{Traefik}, um \en{proxy} reverso e balanceador de carga, isso é, simplifica o roteamento entre vários serviços (microserviços e aplicações em \en{Docker}) \cite{traefik_docs}; \en{Caddy}, plataforma para servir \en{sites}, aplicações e serviços \cite{caddy_docs}; \en{KrakenD}, simplifica acesso à mutiplos serviços de \en{backend}, ele faz isso por centralizar, acelerar e proteger as conexões \cite{krakend_docs}.
